\documentclass[12pt,a4paper]{article}

\usepackage[utf8]{inputenc}
\usepackage[T2A]{fontenc}
\usepackage[ukrainian]{babel}

\usepackage{geometry}
\geometry{left=2cm,right=2cm,top=2cm,bottom=2cm}

\usepackage{graphicx}
\usepackage{array,multirow}
\usepackage{booktabs}
\usepackage{caption}

\usepackage{subcaption}

\usepackage{tcolorbox}

\renewcommand{\thetable}{№\arabic{table}}
\captionsetup[table]{name=Таблиця}

\usepackage{amsmath}
\usepackage{pdflscape}
\usepackage{xcolor}

\usepackage{fvextra}
\usepackage{upquote}

\usepackage{hyperref}

\captionsetup[subfigure]{labelformat=simple,labelsep=space}

\usepackage{minted}

\usemintedstyle{vs}
\setminted{linenos, breaklines, frame=single, fontsize=\small}

\RecustomVerbatimCommand{\VerbatimInput}{VerbatimInput}{
    breaklines=true,
    breaksymbolleft={}
}

% Узагальнений стиль для коду
\newminted{java}{
  fontsize=\small,
  baselinestretch=1.05,
  breaklines,
  linenos,
  frame=lines,
  framesep=2mm
}

% Консоль/PowerShell як plain text
\newminted{text}{
  fontsize=\small,
  breaklines,
  frame=single,
  bgcolor=gray!5
}

\begin{document}

    \begin{titlepage}

        % Картинка шапки (по центру, з відступом від верху)
        \begin{center}
\includegraphics[width=0.7\textwidth]{../../../TITLES/letterhead.png}\end{center}

        \begin{center}
            \large Національний технічний університет України\\
            «Київський політехнічний інститут імені Ігоря Сікорського»\\[1em]
            Факультет інформатики та обчислювальної техніки\\
            Кафедра обчислювальної техніки
        \end{center}

        \vfill

        \begin{center}
            \textbf{\LARGE Інженерія програмного забезпечення}\\[2em]
            \textbf{\Large Практична робота №6}\\[0.5em]
            «Породжувальні шаблони. Шаблони Prototype, Singleton, Factory Method, Abstract Factory, Builder»
        \end{center}

        \vfill

        \begin{flushright}
            Виконав: студент 2 курсу ФІОТ, гр. ІО-41\\
            \textit{Давидчук А.~М.}\\
            Залікова книжка № 4106\\[1em]
            Перевірив: \textit{Васильєва М.~Д.}
        \end{flushright}

        \vfill

        \begin{center}Київ --- 2025\end{center}

    \end{titlepage}

    % Основний документ:

    \setlength{\parindent}{0em}

    \textbf{\underline{Тема:}} «Породжувальні шаблони. Шаблони Prototype, Singleton, Factory Method, Abstract Factory, Builder».

    \vspace{1em}

    \textbf{\underline{Мета:}} Вивчення шаблонів поведінки. Отримання базових навичок з застосування
шаблонів Prototype, Singleton, Factory Method Abstract Factory та Builder.

    \begin{center} \textbf{\Large Практична частина (Завдання №1)} \end{center}

    \setlength{\parindent}{1.5em}

    Мій номер залікової книги: $4106$, звідси $4106 \text{ mod } 7 = 4$. Тому умова до 1 завдання буде:

    Необхідно створити модель для подання розпорядника гри (\texttt{GameManager}), 
    який складається з ігрового простору та списку ігрових фішок. 
    Потрібно забезпечити можливість створення \textbf{тільки одного примірника} 
    розпорядника за допомогою шаблону \textit{Singleton}.

    \subsection*{Вимоги}

    \begin{itemize}
        \item \textbf{Клас розпорядника гри (\texttt{GameManager}):}
        \begin{itemize}
            \item має бути єдиною точкою керування грою;
            \item повинен зберігати стан ігрового простору 
            (наприклад, поле, карта, рівень);
            \item повинен містити список ігрових фішок 
            (об’єкти, які представляють гравців або елементи гри);
            \item має забезпечувати методи для додавання, видалення 
            та оновлення фішок;
            \item має містити метод для запуску або перезапуску гри.
        \end{itemize}

        \item \textbf{Гарантія єдиного примірника:}
        \begin{itemize}
            \item якщо інші частини програми спробують створити новий об’єкт 
            \texttt{GameManager}, це повинно бути неможливо;
            \item усі звернення до розпорядника гри мають здійснюватися 
            лише через уже існуючий примірник.
        \end{itemize}

        \item \textbf{Допоміжні класи (за бажанням для зручності):}
        \begin{itemize}
            \item \texttt{GameSpace} --- описує ігрове поле або середовище;
            \item \texttt{GamePiece} --- представляє окрему ігрову фішку 
            (позиція, тип, стан).
        \end{itemize}

        \item \textbf{Загальна структура системи:}
        \begin{itemize}
            \item система повинна бути спроєктована так, щоб усі елементи гри 
            взаємодіяли між собою \textbf{через} \texttt{GameManager}.
        \end{itemize}
    \end{itemize}

    \subsection*{У методі \texttt{main} продемонструвати:}

    \begin{itemize}
        \item створення кількох «об’єктів» \texttt{GameManager} і показати, 
        що в усіх випадках повертається один і той самий екземпляр 
        (робота шаблону \textit{Singleton});
        \item додавання кількох ігрових фішок у спільний ігровий простір 
        через \texttt{GameManager};
        \item виведення поточного стану гри 
        (наприклад, параметри простору та список фішок);
        \item підтвердити, що зміни, внесені через один доступ 
        до \texttt{GameManager}, відображаються при зверненні до нього 
        з будь-якої іншої частини програми.
    \end{itemize}

    \setlength{\parindent}{0em}

    \textbf{\underline{Загальна структура проекту до компіляції:}}

    \begin{minted}{text}
PS C:\Users\artem\Desktop\reps\university_roadmap\semestr3\software engineering\labs\lab6\task1> tree /f /a
Folder PATH listing
Volume serial number is 1A5B-F8AF
C:.
|   build.xml
|   
\---src
    \---work5
            GameManager.java
            GameManagerDemo.java
            GamePiece.java
            GameSpace.java
    \end{minted}

    \textbf{\underline{Код build.xml задля автоматизації:}}

    \begin{minted}{xml}
<?xml version="1.0" encoding="UTF-8"?>
<project name="task3" default="run" basedir=".">

    <!-- Основнi директорiї проєкту -->
    <property name="src.dir" value="src"/>
    <property name="build.dir" value="build"/>
    <property name="dist.dir" value="dist"/>
    <property name="doc.dir" value="doc"/>

    <!-- Iнiцiалiзацiя: створення службових папок -->
    <target name="init" description="Create build, dist and doc directories">
        <mkdir dir="${build.dir}"/>
        <mkdir dir="${dist.dir}"/>
        <mkdir dir="${doc.dir}"/>
    </target>

    <!-- Компiляцiя вихiдних кодiв -->
    <target name="compile" depends="init" description="Compile source code">
        <javac srcdir="${src.dir}" destdir="${build.dir}" includeantruntime="false" encoding="UTF-8"/>
    </target>

    <!-- Генерацiя Javadoc (документацiя українською) -->
    <target name="javadoc" depends="compile" description="Generate Javadoc documentation">
        <javadoc destdir="${doc.dir}"
                 sourcepath="${src.dir}"
                 packagenames="work5"
                 encoding="UTF-8"
                 charset="UTF-8"
                 docencoding="UTF-8"
                 windowtitle="task3 API"
                 doctitle="task3 API">
        </javadoc>
    </target>

    <!-- Створення jar-файлу -->
    <target name="jar" depends="compile" description="Create JAR archive">
        <jar destfile="${dist.dir}/task3.jar" basedir="${build.dir}">
            <manifest>
                <attribute name="Main-Class" value="work5.GameManagerDemo"/>
            </manifest>
        </jar>
    </target>

    <!-- Запуск демо-програми -->
    <target name="run" depends="compile,javadoc" description="Run demo program">
        <java classname="work5.GameManagerDemo" fork="true">
            <classpath>
                <pathelement path="${build.dir}"/>
            </classpath>
        </java>
    </target>

    <!-- Прибирання результатiв збiрки -->
    <target name="clean" description="Delete build, dist and doc directories">
        <delete dir="${build.dir}"/>
        <delete dir="${dist.dir}"/>
        <delete dir="${doc.dir}"/>
    </target>

</project>
    \end{minted}

    \textbf{\underline{Демонстрація доступних команд збірки:}}

    \begin{minted}{text}
PS C:\Users\artem\Desktop\reps\university_roadmap\semestr3\software engineering\labs\lab6\task1> ant -p
Buildfile: C:\Users\artem\Desktop\reps\university_roadmap\semestr3\software engineering\labs\lab6\task1\build.xml

Main targets:

 clean    Delete build, dist and doc directories
 compile  Compile source code
 init     Create build, dist and doc directories
 jar      Create JAR archive
 javadoc  Generate Javadoc documentation
 run      Run demo program
Default target: run
    \end{minted}

    Вихідний код \,\texttt{*.java}\, буде на віддаленому репозиторії, посилання на яке буде додане разом із звітом.

    \vspace{1em}

    \textbf{\underline{Частина виводу результату збірки демонстрації виконання програми:}}

    \begin{minted}{text}
...

run:
    [java] Checking Singleton:
    [java] manager1 == manager2 -> true
    [java]
    [java] Initial game state (via manager1):
    [java] === Game state ===
    [java] Game number: 1
    [java] Game space: GameSpace{name='Dungeon level', width=20, height=15}
    [java] Pieces:
    [java]   GamePiece{id=1, type='Hero', x=5, y=5}
    [java]   GamePiece{id=2, type='Monster', x=10, y=7}
    [java]   GamePiece{id=3, type='Chest', x=3, y=2}
    [java] ==================
    [java]
    [java] Moving piece with id = 2 (Monster) to (12, 8)...
    [java] Game state after movement (via manager2):
    [java] === Game state ===
    [java] Game number: 1
    [java] Game space: GameSpace{name='Dungeon level', width=20, height=15}
    [java] Pieces:
    [java]   GamePiece{id=1, type='Hero', x=5, y=5}
    [java]   GamePiece{id=2, type='Monster', x=12, y=8}
    [java]   GamePiece{id=3, type='Chest', x=3, y=2}
    [java] ==================
    [java]
    [java] Removing piece with id = 3 (Chest)...
    [java] Game state after removing chest (via manager1):
    [java] === Game state ===
    [java] Game number: 1
    [java] Game space: GameSpace{name='Dungeon level', width=20, height=15}
    [java] Pieces:
    [java]   GamePiece{id=1, type='Hero', x=5, y=5}
    [java]   GamePiece{id=2, type='Monster', x=12, y=8}
    [java] ==================
    [java]
    [java] Restarting game...
    [java] Game state after restart (via manager2):
    [java] === Game state ===
    [java] Game number: 2
    [java] Game space: GameSpace{name='Dungeon level', width=20, height=15}
    [java] Pieces:
    [java]   (no pieces)
    [java] ==================

BUILD SUCCESSFUL
Total time: 2 seconds
    \end{minted}

    \textbf{\underline{Загальна структура проекту після виконання збірки:}}

    \begin{minted}{text}
PS C:\Users\artem\Desktop\reps\university_roadmap\semestr3\software engineering\labs\lab6\task1> tree /f /a
Folder PATH listing
Volume serial number is 1A5B-F8AF
C:.
|   build.xml
|   
+---build
|   \---work5
|           GameManager.class
|           GameManagerDemo.class
|           GamePiece.class
|           GameSpace.class
|
+---dist
+---doc
|   |   allclasses-index.html
|   |   allpackages-index.html
|   |   copy.svg
|   |   element-list
|   |   help-doc.html
|   |   index-all.html
|   |   index.html
|   |   link.svg
|   |   member-search-index.js
|   |   module-search-index.js
|   |   overview-tree.html
|   |   package-search-index.js
|   |   script.js
|   |   search-page.js
|   |   search.html
|   |   search.js
|   |   stylesheet.css
|   |   tag-search-index.js
|   |   type-search-index.js
|   |
|   +---legal
|   |       COPYRIGHT
|   |       jquery.md
|   |       jqueryUI.md
|   |       LICENSE
|   |
|   +---resources
|   |       glass.png
|   |       x.png
|   |
|   +---script-dir
|   |       jquery-3.7.1.min.js
|   |       jquery-ui.min.css
|   |       jquery-ui.min.js
|   |
|   \---work5
|           GameManager.html
|           GameManagerDemo.html
|           GamePiece.html
|           GameSpace.html
|           package-summary.html
|           package-tree.html
|
\---src
    \---work5
            GameManager.java
            GameManagerDemo.java
            GamePiece.java
            GameSpace.java
    \end{minted}

    \textbf{\underline{Скріншот наявності документації:}}

    \begin{figure}[h!]
        \centering
        \includegraphics[width=\linewidth]{1_1.png}
        \caption{Скріншот документації Javadoc для 1 завдання}
    \end{figure}

    \textbf{\underline{UML-діаграма проекту:}}

    \begin{figure}[h!]
        \centering
        \includegraphics[width=0.6\linewidth]{1_2.png}
        \caption{UML-діаграма проекту для 1 завдання}
    \end{figure}

    \newpage

    \begin{center} \textbf{\Large Практична частина (Завдання №2)} \end{center}

    Мій номер залікової книги: $4106$, звідси $4106 \text{ mod } 7 = 4$. Тому умова до 2 завдання буде:

    Необхідно визначити специфікації класів для подання сімейства інструментів
    універсального інтегрованого середовища розробки
    (\texttt{Validator}, \texttt{Compiler}, \texttt{Debugger}) з їх реалізацією
    для різних мов (\texttt{Java}, \texttt{C++}).

    \subsection*{Вимоги}

    \begin{itemize}
        \item \textbf{Фабрика інструментів середовища розробки:}
        \begin{itemize}
            \item створити абстрактний клас або інтерфейс \texttt{IDEToolsFactory},
            який визначає фабричні методи для створення компонентів середовища
            розробки:

            \texttt{createValidator()}~-- повертає об’єкт для перевірки коду,
            \texttt{createCompiler()}~-- повертає об’єкт компілятора,
            \texttt{createDebugger()}~-- повертає об’єкт дебагера;
            \item класи \texttt{JavaIDEToolsFactory} та \texttt{CppIDEToolsFactory}
            є конкретними реалізаціями \texttt{IDEToolsFactory}, що створюють
            інструменти для відповідних мов програмування.
        \end{itemize}

        \item \textbf{Інтерфейси інструментів:}
        \begin{itemize}
            \item створити абстрактні інтерфейси \texttt{Validator},
            \texttt{Compiler} та \texttt{Debugger}, які визначають базові методи роботи,
            наприклад:
            \texttt{validateCode(String code)} для \texttt{Validator},
            \texttt{compile(String source)} для \texttt{Compiler},
            \texttt{debug(String program)} для \texttt{Debugger}.
        \end{itemize}

        \item \textbf{Реалізації для мов програмування:}
        \begin{itemize}
            \item класи \texttt{JavaValidator}, \texttt{JavaCompiler},
            \texttt{JavaDebugger}~-- реалізації інтерфейсів для мови Java;
            \item класи \texttt{CppValidator}, \texttt{CppCompiler},
            \texttt{CppDebugger}~-- реалізації інтерфейсів для мови C++.
        \end{itemize}

        \item \textbf{Розширюваність системи:}
        \begin{itemize}
            \item забезпечити можливість прозорого для клієнта розширення системи
            новими реалізаціями для інших мов без зміни клієнтського коду.
        \end{itemize}
    \end{itemize}

    \subsection*{У методі \texttt{main} продемонструвати:}

    \begin{itemize}
        \item створення фабрики певного типу (наприклад,
        \texttt{JavaIDEToolsFactory} або \texttt{CppIDEToolsFactory});
        \item за допомогою фабрики створити об’єкти всіх компонентів:
        \texttt{Validator}, \texttt{Compiler} і \texttt{Debugger};
        \item виклик відповідних методів для перевірки, компіляції
        та відладки умовного фрагмента коду.
    \end{itemize}

    \textbf{\underline{Загальна структура проекту до компіляції:}}

    \begin{minted}{text}
PS C:\Users\artem\Desktop\reps\university_roadmap\semestr3\software engineering\labs\lab6\task2> tree /f /a
Folder PATH listing
Volume serial number is 1A5B-F8AF
C:.
|   build.xml
|   
\---src
    \---work5
            Compiler.java
            CppCompiler.java
            CppDebugger.java
            CppIDEToolsFactory.java
            CppValidator.java
            Debugger.java
            IDEToolsDemo.java
            IDEToolsFactory.java
            JavaCompiler.java
            JavaDebugger.java
            JavaIDEToolsFactory.java
            JavaValidator.java
            Validator.java
    \end{minted}

    \textbf{\underline{Код build.xml задля автоматизації:}}

    \begin{minted}{xml}
<?xml version="1.0" encoding="UTF-8"?>
<project name="task2" default="run" basedir=".">

    <!-- Основнi директорiї проєкту -->
    <property name="src.dir" value="src"/>
    <property name="build.dir" value="build"/>
    <property name="dist.dir" value="dist"/>
    <property name="doc.dir" value="doc"/>

    <!-- Iнiцiалiзацiя: створення службових папок -->
    <target name="init" description="Create build, dist and doc directories">
        <mkdir dir="${build.dir}"/>
        <mkdir dir="${dist.dir}"/>
        <mkdir dir="${doc.dir}"/>
    </target>

    <!-- Компiляцiя вихiдних кодiв -->
    <target name="compile" depends="init" description="Compile source code">
        <javac srcdir="${src.dir}"
               destdir="${build.dir}"
               includeantruntime="false"
               encoding="UTF-8"/>
    </target>

    <!-- Генерацiя Javadoc (українська документацiя) -->
    <target name="javadoc" depends="compile" description="Generate Javadoc documentation">
        <javadoc destdir="${doc.dir}"
                 sourcepath="${src.dir}"
                 packagenames="work5"
                 encoding="UTF-8"
                 charset="UTF-8"
                 docencoding="UTF-8"
                 windowtitle="IDE Tools Factory API"
                 doctitle="IDE Tools Factory API">
        </javadoc>
    </target>

    <!-- Створення jar-файлу -->
    <target name="jar" depends="compile" description="Create JAR archive">
        <jar destfile="${dist.dir}/ide-tools-factory.jar" basedir="${build.dir}">
            <manifest>
                <attribute name="Main-Class" value="work5.IDEToolsDemo"/>
            </manifest>
        </jar>
    </target>

    <!-- Запуск демо-програми -->
    <target name="run" depends="compile,javadoc" description="Run demo program">
        <java classname="work5.IDEToolsDemo" fork="true">
            <classpath>
                <pathelement path="${build.dir}"/>
            </classpath>
        </java>
    </target>

    <!-- Очищення результатiв збiрки -->
    <target name="clean" description="Delete build, dist and doc directories">
        <delete dir="${build.dir}"/>
        <delete dir="${dist.dir}"/>
        <delete dir="${doc.dir}"/>
    </target>

</project>
    \end{minted}

    \textbf{\underline{Демонстрація доступних команд збірки:}}

    \begin{minted}{text}
PS C:\Users\artem\Desktop\reps\university_roadmap\semestr3\software engineering\labs\lab6\task2> ant -p 
Buildfile: C:\Users\artem\Desktop\reps\university_roadmap\semestr3\software engineering\labs\lab6\task2\build.xml

Main targets:

 clean    Delete build, dist and doc directories
 compile  Compile source code
 init     Create build, dist and doc directories
 jar      Create JAR archive
 javadoc  Generate Javadoc documentation
 run      Run demo program
Default target: run
    \end{minted}

    Вихідний код \,\texttt{*.java}\, буде на віддаленому репозиторії, посилання на яке буде додане разом із звітом.

    \vspace{1em}

    \textbf{\underline{Частина виводу результату збірки демонстрації виконання програми:}}

    \begin{minted}{text}
run:
     [java] Using Java IDE tools factory:
     [java] JavaValidator: validating Java code...
     [java]   Code: System.out.println("Hello, Java!");
     [java]   Result: no syntax errors found.
     [java] JavaCompiler: compiling Java source...
     [java]   Source: System.out.println("Hello, Java!");
     [java]   Result: compilation successful, .class file generated.
     [java] JavaDebugger: starting Java debug session...
     [java]   Program: JavaSampleProgram
     [java]   Result: breakpoint hit at line 42.
     [java]
     [java] Using C++ IDE tools factory:
     [java] CppValidator: validating C++ code...
     [java]   Code: std::cout << "Hello, C++!" << std::endl;
     [java]   Result: no syntax errors found.
     [java] CppCompiler: compiling C++ source...
     [java]   Source: std::cout << "Hello, C++!" << std::endl;
     [java]   Result: compilation successful, executable generated.
     [java] CppDebugger: starting C++ debug session...
     [java]   Program: CppSampleProgram
     [java]   Result: watchpoint triggered on variable x.

BUILD SUCCESSFUL
Total time: 1 second
    \end{minted}

    \textbf{\underline{Загальна структура проекту після виконання збірки:}}

    \begin{minted}{text}
PS C:\Users\artem\Desktop\reps\university_roadmap\semestr3\software engineering\labs\lab6\task2> tree /f /a
Folder PATH listing
Volume serial number is 1A5B-F8AF
C:.
|   build.xml
|   
+---build
|   \---work5
|           Compiler.class
|           CppCompiler.class
|           CppDebugger.class
|           CppIDEToolsFactory.class
|           CppValidator.class
|           Debugger.class
|           IDEToolsDemo.class
|           IDEToolsFactory.class
|           JavaCompiler.class
|           JavaDebugger.class
|           JavaIDEToolsFactory.class
|           JavaValidator.class
|           Validator.class
|
+---dist
+---doc
|   |   allclasses-index.html
|   |   allpackages-index.html
|   |   copy.svg
|   |   element-list
|   |   help-doc.html
|   |   index-all.html
|   |   index.html
|   |   link.svg
|   |   member-search-index.js
|   |   module-search-index.js
|   |   overview-tree.html
|   |   package-search-index.js
|   |   script.js
|   |   search-page.js
|   |   search.html
|   |   search.js
|   |   stylesheet.css
|   |   tag-search-index.js
|   |   type-search-index.js
|   |
|   +---legal
|   |       COPYRIGHT
|   |       jquery.md
|   |       jqueryUI.md
|   |       LICENSE
|   |
|   +---resources
|   |       glass.png
|   |       x.png
|   |
|   +---script-dir
|   |       jquery-3.7.1.min.js
|   |       jquery-ui.min.css
|   |       jquery-ui.min.js
|   |
|   \---work5
|           Compiler.html
|           CppCompiler.html
|           CppDebugger.html
|           CppIDEToolsFactory.html
|           CppValidator.html
|           Debugger.html
|           IDEToolsDemo.html
|           IDEToolsFactory.html
|           JavaCompiler.html
|           JavaDebugger.html
|           JavaIDEToolsFactory.html
|           JavaValidator.html
|           package-summary.html
|           package-tree.html
|           Validator.html
|
\---src
    \---work5
            Compiler.java
            CppCompiler.java
            CppDebugger.java
            CppIDEToolsFactory.java
            CppValidator.java
            Debugger.java
            IDEToolsDemo.java
            IDEToolsFactory.java
            JavaCompiler.java
            JavaDebugger.java
            JavaIDEToolsFactory.java
            JavaValidator.java
            Validator.java
    \end{minted}

    \textbf{\underline{Скріншот наявності документації:}}

    \begin{figure}[h!]
        \centering
        \includegraphics[width=\linewidth]{2_1.png}
        \caption{Скріншот документації Javadoc для 2 завдання}
    \end{figure}

    \textbf{\underline{UML-діаграма проекту:}}

    \begin{figure}[h!]
        \centering
        \includegraphics[width=1.0\linewidth]{2_2.png}
        \caption{UML-діаграма проекту для 2 завдання}
    \end{figure}

    \newpage

    \setlength{\parindent}{1.5em}

    \noindent
    \textbf{\underline{Висновок:}} У ході виконання роботи було спроєктовано та реалізовано модель
сімейства інструментів інтегрованого середовища розробки (IDE),
що включає компоненти \texttt{Validator}, \texttt{Compiler} та
\texttt{Debugger} для двох мов програмування --- Java та C++.
Для створення взаємопов'язаних об'єктів цих типів було застосовано
шаблон проєктування \textit{Abstract Factory} у вигляді інтерфейсу
\texttt{IDEToolsFactory} та його конкретних реалізацій
\texttt{JavaIDEToolsFactory} і \texttt{CppIDEToolsFactory}.
Клієнтський код (\texttt{IDEToolsDemo}) працює виключно з абстракціями
фабрики та інтерфейсами інструментів, що демонструє ослаблене
зв'язування між використанням інструментів і їх конкретними реалізаціями.

Результати запуску програми показали коректність роботи фабрик:
для кожного типу фабрики створюється узгоджений набір інструментів
(валідатор, компілятор, налагоджувач), над якими послідовно виконуються
операції перевірки коду, компіляції та налагодження. Завдяки такій
організації системи стає можливим прозоре розширення: для додавання
підтримки нової мови достатньо реалізувати нові класи-продукти
(\texttt{XxxValidator}, \texttt{XxxCompiler}, \texttt{XxxDebugger})
та відповідну фабрику \texttt{XxxIDEToolsFactory}, не змінюючи клієнтський код.

Таким чином, мету роботи досягнуто: розроблено узагальнену структуру
IDE-інструментів із використанням шаблону \textit{Abstract Factory},
продемонстровано переваги інкапсуляції способів створення об'єктів
і забезпечено розширюваність програмної системи без модифікації
існуючого клієнтського коду.

\end{document}